\section{Experiments}
\label{sec:experiments}
In this section, we focus on the empirical evaluation of our desiderata.
We compare Ensembles, Monte Carlo dropout (MCD), evidential regression, SVGP-DKL\footnote{For SVGP-DKL versus LNK, see \autoref{app:svgp_vs_lnk}}, and \oursacro{}.
We pair each of the UE methods with SchNet~\citep{schutt2021equivariant}, DimeNet++\citep{gasteigerDirectionalMessagePassing2019,gasteigerFastUncertaintyAwareDirectional2020} and NequIP~\citep{batznerEquivariantGraphNeural2022}.
As traditional GPs and SVGP either do not allow for varying input sizes or result in poor predictive performance, we do not discuss these here.
Further, as \desone{} and \destwo{} are fixed desiderata, discussed in earlier \autoref{sec:survey} and \autoref{sec:method}, we do not perform experiments for those.

In our first experiment, we analyze \desfour{} by testing whether the UE is size consistent.
Next, we analyze the \desfive{} on two common QM datasets, MD17~\citep{chmiela2017machine} and QM7X~\citep{hojaQM7XComprehensiveDataset2021}.
To measure \dessix{}, we look at the calibration score and out-of-distribution (OOD) detection.
For calibration, we use the calibration regression score from \citet{charpentier_natural_2022}.
For OOD detection we use the area under the receiver operating characteristic curve (AUC-ROC). For each backbone, we use the respective hyperparameters from their papers, listed in \autoref{app:baselinedetails}.
If not further specified, we train each model on the standard combined energy and force loss~\citep{gasteigerDirectionalMessagePassing2019}: 
\begin{equation}
\mathcal{L}_\text{total} = (1-\rho)\mathcal{L}_\text{energy} + \rho\mathcal{L}_\text{force}.
\end{equation}
Here $\rho$ is a hyperparameter to set the importance of energy and force loss. For $\text{\energy}, \text{force}$, we specify the loss in \cref{app:baselinedetails}.

\textbf{Methods.}
\underline{\emph{Ensembles:}} For Ensembles we use a sample size of $5$ differently initialized models trained on the same dataset.
\underline{\emph{Monte-Carlo Dropout:}} Similar to Ensembles, we choose a sample size of 5 for evaluating uncertainties. 
As dropout rate, we use a rate of \SI{20}{\percent}.
We study different dropout rates in \autoref{app:dropout_exp} but found \SI{20}{\percent} to be good tradeoff between \desfive{} amd \dessix{}.
\underline{\emph{Evidential regression:}} We adapt the energy loss according to \cite{soleimany2021evidential} instead of the MAE. Note that we use a smaller coefficient for the forces loss of $\rho=0.9$ (instead of $\rho=0.99$ used with other models).
In early experiments, we found larger force coefficients leading to poor uncertainty estimates.
\underline{\emph{SVGP-DKL and \oursacro{}:}} As described in \autoref{sec:method}, we use the Predictive Log-Likelihood as energy loss instead of the MAE. We compare \oursacro{} against SVGP-DKL in \cref{app:svgp_vs_lnk} to highlight the importance of localization. \underline{Graph Mixture Density Networks:} we investigate the performance of this approach in \cref{app:gmdn}.

\textbf{Datasets.}
\underline{\emph{QM7-X:}} \cite{hojaQM7XComprehensiveDataset2021} This dataset covers both equilibrium and non-equilibrium structures. We train on equilibrium structures and non-equilibrium structures are considered OOD data. \underline{\emph{MD17:}} \cite{chmiela2017machine} MD17 contains energies and forces for molecular dynamics trajectories of different organic molecules. %

\subsection{\desfour{}}\label{sec:exp_extensivity}
\begin{figure}
    \centering
    \includegraphics[width=\linewidth]{figures/size_consistency_experiment.pdf}
    \caption{$u_{\energy}$ for increasing number of identical molecules with sufficient distance to one another.}
    \label{fig:extensivity}
\end{figure}
In this experiment, we analyze the \desfour{} property for all UE models.
For this, we duplicate the same molecule translated by \SI{15}{\angstrom}\footnote{The distance larger than the cutoff of all GNNs.}.
We measure the total uncertainty of the increasing number of molecules.
For a model fulfilling \desfour{}, the total uncertainty should be proportional to the number of molecules, i.e., $\uncertainty_\energy \propto i$ where $i$ is the number of copies.

\autoref{fig:extensivity} depicts the dependency of the uncertainty by the number of molecules.
We observe that both SVGP-DKL and Evidential do not exhibit the desired behavior. The increase of uncertainty is too steep and after three molecules, stops increasing. Ensembles and \oursacro{} meet the requirements of increasing linearly with slope $1$. We can observe that MCD exhibits an increase in uncertainty with the growing number of molecules but due to its randomness, \desfour{} will only be satisfied in expectation. 

\subsection{\desfive{}}
\begin{table}
\caption{MAE of SchNet on MD17 with different UE methods (energies in $\si{\kcal\per\mole}$, forces in $\si{\kcal\per\mole\per\angstrom}$)}\label{tab:md17_schnet}
\resizebox{\columnwidth}{!}{
\begin{tabular}{llccc||c|c}
\toprule
{} & & MCD & Evidential & fixed-\oursacro{} & Backbone & Ensemble \\
\midrule
\multirow{3}{*}{aspirin} & Energy & 2.336 & 1.007 & \textbf{0.325} & 0.321 & 0.272\\ & Forces & 2.055 & 1.088 & \textbf{0.865} & 0.852 & 0.701\\ & Calibration & 1.256 & \textbf{1.077} & 1.370 & - & 1.354\\
\hline
\multirow{3}{*}{ethanol} & Energy & 0.646 & 0.081 & \textbf{0.065} & 0.064 & 0.06\\ & Forces & 0.925 & 0.328 & \textbf{0.271} & 0.112 & 0.191\\ & Calibration & \textbf{1.288} & 1.643 & 1.566 & - & 1.469\\
\hline
\multirow{3}{*}{MDA} & Energy & 1.493 & 0.130 & \textbf{0.111} & 0.111 & 0.098\\ & Forces & 1.467 & 0.535 & \textbf{0.481} & 0.461 & 0.368\\ & Calibration & \textbf{1.257} & 1.580 & 1.475 & - & 1.442\\
\hline
\multirow{3}{*}{Nap} & Energy & 2.739 & 0.346 & \textbf{0.129} & 0.131 & 0.124\\ & Forces & 1.385 & 0.357 & \textbf{0.266} & 0.255 & 0.211\\ & Calibration & 1.280 & \textbf{1.258} & 1.427 & - & 1.510\\
\hline
\multirow{3}{*}{SAC} & Energy & 2.301 & 0.165 & \textbf{0.142} & 0.141 & 0.132\\ & Forces & 1.853 & 0.607  & \textbf{0.481} & 0.465 & 0.378\\ & Calibration & \textbf{1.266} & 1.527 & 1.504 & - & 1.483\\
\hline
\multirow{3}{*}{toluene} & Energy & 1.662 & 0.202 & \textbf{0.106} & 0.108 & 0.099\\ & Forces & 1.404 & 0.402 & \textbf{0.310} & 0.296 & 0.230\\ & Calibration & \textbf{1.269} & 1.472 & 1.501 & - & 1.526\\
\hline
\multirow{3}{*}{uracil} & Energy & 1.759 & 0.168 & \textbf{0.121} & 0.119 & 0.115\\ & Forces & 1.945 & 0.480 & \textbf{0.338} & 0.326 & 0.265\\ & Calibration & \textbf{1.262} & 1.524 & 1.440 & - & 1.543\\
\bottomrule
\end{tabular}
}
\end{table}
For this task, we use the fixed-\oursacro{} as it enables us to reuse one ensemble backbone while showing the comparison of how well the pure GNN can perform on that dataset. Hence, we also list the backbone performance as a reference. 
We compare the average predictive performance across all seven MD17 molecules for each combination of backbone and UE method in \autoref{tab:md17_average}.
One may see that while MCD and evidential regression sacrifice accuracy, \oursacro{} maintains backbone-like predictive performance. 
For SchNet, we provide per-molecule results in \autoref{tab:md17_schnet}. We show additional results for other backbones in \autoref{app:additional_experiments}.

\begin{table}
    \centering
    \caption{Average metrics over seven molecules of MD17 for different backbones}
    \vspace{0.05in}
    \label{tab:md17_average}
    \resizebox{\columnwidth}{!}{
    \begin{tabular}{llccc|c|c}
\toprule
   {} & {} & Dropout & Evidential & fixed-\oursacro{} & Backbone & Ensenmble \\
   \midrule
\multirow{2}{*}{DimeNet++} & Energy & 1.126 & 0.104 & \textbf{0.106} & 0.109 & 0.144\\ & Forces & 0.624 & 0.2 & \textbf{0.196} & 0.217 & 0.219\\ %
\hline
\multirow{2}{*}{SchNet} & Energy & 1.848 & 0.3 & \textbf{0.13} & 0.142 & 0.129\\ & Forces & 1.576 & 0.542 & 0.423 & 0.395 & 0.335\\% & Calibration & \\
\hline
\multirow{2}{*}{NequIP} & Energy & 1.557 & \textbf{0.098} & 0.102 & 0.101 & 0.098\\ & Forces & 0.655 & 0.405 & \textbf{0.342} & 0.147 & 0.106\\% & Calibration & \\

\bottomrule
\end{tabular}
    }
\end{table}

In \autoref{tab:qm7x_dimenet_predictive}, we can see the predictive performance on QM7X.
In energy prediction, we again observe that Ensembles, Evidential, and \oursacro{} perform well in terms of accuracy.
Remarkably, MCD's energy accuracy significantly deteriorates with increasing dropout rates.
For forces, Ensembles, and \oursacro{} perform well. But, unlike energies, the Evidential approach significantly worsens predictive results.
Despite its worse energy error, MCD achieves low force errors.

\begin{table}
    \caption{MAE of Dimenet++ trained on QM7X-Equilibrium with different UE methods (energies in $\si{\eV}$, forces in $\si{\eV\per\angstrom}$)}
    \vspace{0.05in}
    \resizebox{\columnwidth}{!}{
    \begin{tabular}{lccccc|c}
\toprule
{} & MCD 20\% & MCD 1\% & Evidential & fixed-\oursacro{} & \oursacro{} & Ensemble\\
\midrule
Energy  &   0.24 & 0.09 & \textbf{0.02} & 0.046 & \textbf{0.0208} & 0.0523\\
Forces  &   0.0087 & \textbf{0.0043} & 0.0225 & 0.047 & 0.0046 & 0.0034\\
\midrule
Calibration    &  \textbf{1.279} & 1.285 & 1.629 & 1.293 & 1.339 & 2.597 \\
\bottomrule
\end{tabular}


    }
    \label{tab:qm7x_dimenet_predictive}
\end{table}

\subsection{\dessix{}}

Now, we measure the uncertainty estimates for different approaches. For SchNet on MD17, \autoref{tab:md17_schnet} lists the calibration scores for each UE method on each molecule.
Here, most approaches perform similarly while MCD surprisingly yields the the best results.
We observe a similar pattern for QM7X in \autoref{tab:qm7x_dimenet_predictive}. However, the energy prediction of MCD with $0.24$ is more than a magnitude worse than others and as a uniform distribution minimizes the calibration score, it only has significance if the energy prediction is good. All other methods have similar calibration scores.

\begin{figure}
\centering
\includegraphics[width=1.0\columnwidth]{figures/NEW_schnet_md17_all_methods_2x2.pdf}
\caption{Heatmap of AUC-ROC values from SchNet on MD17. Each row corresponds to a separate model trained on the molecule written on the left and tested on all other molecules.}\label{fig:md17_uq}
\end{figure}

In the following we test the performance of the methods on left-out class detection: we train a model on one of seven molecules and see whether it can distinguish the others based on its uncertainty estimate. An ideal estimator will have small values for the trained molecule and higher values for the other ones. This would enable a potential user to set a threshold to trust the model prediction whenever its uncertainty is lower than that threshold. To measure performance, we compute the area under the receiver operator curve of the uncertainty scores for ID and OOD data. A value closer to 1 indicates a better separation of the uncertainty scores, which means the model is better at detecting when it is operating on uncertain data.

For Evidential and \oursacro{} we use a single forward pass uncertainty estimator and for MCD and Ensembles we can use the variance of the energy prediction. In \autoref{fig:md17_uq}, we can see heat maps depicting the AUC-ROC score for each pairwise combination with SchNet as backbone. The rows show the molecule that the model is trained on and the off-diagonal columns are the respective OOD sample. On the diagonal, we expect a score of 0.5 while a score of 1 is optimal on the off-diagonals.

We observe that MCD's and Evidential's energy uncertainties are insufficient for such a task as both results exhibit no clear diagonal structure. Meanwhile, Ensembles exhibit perfect separation between in and out-of-distribution samples. With significantly fewer computational resources, \oursacro{} achieves identical results to Ensembles in a single forward pass, see \autoref{app:additional_experiments} for further results. This shows that both approaches can distinguish the molecule distribtions perfectly. 

We further evaluate the uncertainty estimation on QM7X, where we use the equilibrium structures to train our model and try to distinguish the out-of-equilibrium (OOE) ones. This task is more challenging as the distribution of the OOD samples is closer to the ID ones as OOE structures are generated by perturbing ID molecules. 

In \autoref{tab:qm7x_uq}, we can see the AUC-ROC scores with DimeNet++ as backbone. For MCD, we tested two dropout rates and neither could distinguish OOD samples. We observe similar behavior for the Evidential model.
In contrast, Ensembles again achieve close to perfect separation.
\oursacro{} is located in between these two extremes with an AUC-ROC score of 0.803, i.e., 2.1 times lower error than Evidential. Note, that this is obtained with a forward pass at inference, compared to five from the Ensemble. 
\begin{table}
    \centering
    \caption{OOD detection on QM7X based on the energy prediction with DimeNet$++$ backbone.}
    \vspace{0.05in}
    \resizebox{\columnwidth}{!}{
    \begin{tabular}{lccccc|c}
\toprule
{} & MCD 20\% & MCD 1\% & Evidential & fixed-\oursacro{} & \oursacro{} & Ensemble\\
\midrule
AUC-ROC & 0.511 & 0.504 & 0.587 & 0.742 & \textbf{0.803} & 0.947\\
\bottomrule
\end{tabular}


    }
    \label{tab:qm7x_uq}
\end{table}
\subsection{\desseven{}}
Evidential and \oursacro{} both only require one single forward pass\footnote{With the backpropagation to calculate the forces} while MCD and Ensembles use 5 in our experiments. Thus it is natural that these techniques offer a higher \desseven{}. We evaluate the average time per sample for DimeNet++ on QM7X in \autoref{tab:speed}. We can see that both Evidential and \oursacro{} are indeed  more than $4.6\times$ faster than the variance based approaches. We further observe that \oursacro{} and Evidential are having very similar runtime which is due to similar number of parameters, as we in \cref{app:complexity}. Lastly, Ensembles are not only the having the computational demand at inference, but we also have to train multiple models which is additional overhead, especially in an active learning setting where one has to retrain the model many times.
\begin{table}
    \centering
    \caption{Inference runtime in $\si\ms$ on QM7X data for each UE method.}
    \vspace{0.05in}
    \resizebox{\columnwidth}{!}{
    \begin{tabular}{lcccc}
\toprule
{} & \oursacro{} & Evidential & MCD & Ensemble\\
\midrule
Runtime [ms]  & $59.3 \pm 2.68$ & $58.8 \pm 5.35$ & $229 \pm 0.25$ & $274 \pm 8.25$\\
\#Trainings         & 1 & 1 & 1 & 5\\
\bottomrule
\end{tabular}
    }
    \label{tab:speed}
\end{table}
\subsection{Force Uncertainty}
As one cannot obtain force uncertainties with a single forward pass~\citep{gasteigerFastUncertaintyAwareDirectional2020}, force uncertainties directly violate our \desseven{} desiderata.
Nontheless, to provide a broad overview we here look at force uncertainties for OOE detection.
As single-forward pass methods cannot model $u_{\forces}$, we use the MCD approach on LNK by leveraging the dropout at inference time, denoted as MCD-\oursacro{}. We compare against Ensembles and MCD. For all methods we obtain the force uncertainties by computing the trace of the empirical covariance matrix.

\autoref{tab:force} lists the predictive performances and AUC-ROC scores for different dropout rates.
We observe that MCD-\oursacro{} consistently outperforms MCD both in predictive performance as well as in OOE detection.
Further, as the error increases significantly for MCD with increasing dropout rates, we find MCD-\oursacro{} to scale stably with clear improvements in OOE detection.
While Ensembles report close to perfect separation, it is also the computationally most intensive method due to its multiple trainings.

\begin{table}[htb]
    \centering
    \caption{Comparison of Ensembles, MCD, and \oursacro{} on QM7X with DimeNet++ as backbone for different dropout rates.}
    \vspace{0.05in}
    \resizebox{\columnwidth}{!}{
    \begin{tabular}{llccccc}
\toprule
{Method} &  {Metric} & 0\% & 1\% & 5\% & 10\% & 20\%  \\
\midrule
\multirow{4}{*}{MCD} & Energy MAE & -- & 0.09 & 0.134 & 0.259 & 0.24  \\
& Forces MAE & -- & 0.004 & 0.005 & 0.012 & 0.009  \\
& AUC-ROC energy & -- &  0.504 & 0.511 & 0.55 & 0.511 \\
& AUC-ROC forces & -- &  0.86 & 0.87 & 0.89 & 0.91  \\
\midrule
\multirow{4}{*}{MCD-\oursacro{}} & Energy MAE & -- & 0.059 & 0.024 & 0.039 & 0.036\\
& Forces MAE & -- & 0.007 & 0.004 & 0.006 & 0.007\\
& AUC-ROC energy & -- & 0.6 & 0.69 & 0.72 & 0.72\\
& AUC-ROC forces & -- & 0.79 & 0.82 & 0.84 & 0.9\\
\midrule
\multirow{4}{*}{Ensemble} & Energy MAE & 0.052 & -- & -- & -- & -- \\
& Forces MAE & 0.003  & -- & -- & -- & --\\
& AUC-ROC energy & 0.947 & -- & -- & -- & --\\
& AUC-ROC forces & 0.996 & -- & -- & -- & -- \\
\bottomrule
\end{tabular}
    }
    \label{tab:force}
\end{table}